\documentclass[lualatex,ja=standard,a4paper,10pt]{bxjsarticle}

\geometry{margin=17mm,headheight=14pt}
\usepackage{luatexja-fontspec}
\setmainjfont{Noto Serif CJK JP}
\setsansjfont{Noto Sans CJK JP}
\setmonojfont{Noto Sans Mono CJK JP}
\usepackage{graphicx}
\usepackage{booktabs}
\usepackage{tabularx}
\usepackage{array}
\usepackage{amsmath}
\usepackage{xcolor}
\usepackage{caption}
\usepackage{subcaption}
\usepackage{float}
\usepackage{enumitem}
\usepackage{fancyhdr}
\usepackage{xurl}
\usepackage{titlesec}
\usepackage[
  colorlinks=true,
  linkcolor=blue!45!black,
  citecolor=blue!45!black,
  urlcolor=blue!55!black,
  pdfauthor={杉中史哉},
  pdftitle={YOLO-WorldのFailure Case Analysisと改善}
]{hyperref}

\graphicspath{
  {results/yoloworld/figures/}
  {results/yoloworld/examples/}
}
% Generated from experiment CSV/JSON files. Do not edit by hand.
\providecommand{\TotalImageCount}{100}
\providecommand{\ValidationImageCount}{30}
\providecommand{\TestImageCount}{70}
\providecommand{\ValidationGTCount}{58}
\providecommand{\TestGTCount}{128}
\providecommand{\ValidationCategoryGTMin}{11}
\providecommand{\ValidationCategoryGTMax}{18}
\providecommand{\TestSmallGTCount}{21}
\providecommand{\ManifestPath}{data/processed/coco100\_manifest.json}
\providecommand{\GammaValidationTieCount}{1}
\providecommand{\SubsetMaxVariantCount}{2}
\providecommand{\DevelopmentImageCount}{200}
\providecommand{\SameModelDevelopmentImageCount}{300}
\providecommand{\SameModelHoldoutImageCount}{100}
\providecommand{\SameModelHoldoutGTCount}{186}
\providecommand{\SameModelName}{yolov8s-worldv2.pt}
\providecommand{\SameModelBaselineCleanFOne}{0.638}
\providecommand{\SameModelImprovedCleanFOne}{0.634}
\providecommand{\SameModelCleanFOneDelta}{-0.003}
\providecommand{\SameModelBaselineBlurTwoFOne}{0.589}
\providecommand{\SameModelImprovedBlurTwoFOne}{0.590}
\providecommand{\SameModelBaselineBlurTwoAP}{0.510}
\providecommand{\SameModelImprovedBlurTwoAP}{0.497}
\providecommand{\SameModelBaselineBlurFourFOne}{0.339}
\providecommand{\SameModelImprovedBlurFourFOne}{0.505}
\providecommand{\SameModelBaselineBlurFourAP}{0.256}
\providecommand{\SameModelImprovedBlurFourAP}{0.404}
\providecommand{\SameModelBaselineBlurFourRecall}{0.226}
\providecommand{\SameModelImprovedBlurFourRecall}{0.382}
\providecommand{\SameModelBaselineBlurMeanFOne}{0.464}
\providecommand{\SameModelImprovedBlurMeanFOne}{0.548}
\providecommand{\SameModelBlurMeanFOneDelta}{+0.084}
\providecommand{\SameModelBlurMeanFOneDeltaLow}{+0.047}
\providecommand{\SameModelBlurMeanFOneDeltaHigh}{+0.122}
\providecommand{\SameModelDevelopmentFOneDelta}{+0.076}
\providecommand{\SameModelDevelopmentFOneDeltaLow}{+0.052}
\providecommand{\SameModelDevelopmentFOneDeltaHigh}{+0.099}
\providecommand{\SameModelMinimumFoldGain}{+0.061}
\providecommand{\WienerRegularization}{0.001}
\providecommand{\WienerBlend}{0.75}
\providecommand{\BlurCleanThreshold}{36.68}
\providecommand{\BlurSeverityThreshold}{3.35}
\providecommand{\SameModelDetectorCorrect}{295}
\providecommand{\SameModelDetectorTotal}{300}
\providecommand{\SameModelCleanCorrect}{99}
\providecommand{\PromptPrototypeDevelopmentImageCount}{400}
\providecommand{\PromptPrototypeHoldoutImageCount}{100}
\providecommand{\PromptPrototypeHoldoutGTCount}{228}
\providecommand{\PromptPrototypeAnchorWeight}{0.75}
\providecommand{\PromptPrototypeBaselineCanonicalFOne}{0.605}
\providecommand{\PromptPrototypeImprovedCanonicalFOne}{0.605}
\providecommand{\PromptPrototypeBaselineSynonymFOne}{0.389}
\providecommand{\PromptPrototypeImprovedSynonymFOne}{0.583}
\providecommand{\PromptPrototypeBaselineHypernymFOne}{0.372}
\providecommand{\PromptPrototypeImprovedHypernymFOne}{0.566}
\providecommand{\PromptPrototypeBaselineTargetMeanFOne}{0.380}
\providecommand{\PromptPrototypeImprovedTargetMeanFOne}{0.574}
\providecommand{\PromptPrototypeTargetMeanFOneDelta}{+0.194}
\providecommand{\PromptPrototypeTargetMeanFOneDeltaLow}{+0.135}
\providecommand{\PromptPrototypeTargetMeanFOneDeltaHigh}{+0.254}
\providecommand{\PromptPrototypeBaselineTargetMeanAP}{0.443}
\providecommand{\PromptPrototypeImprovedTargetMeanAP}{0.511}
\providecommand{\PromptPrototypeTargetMeanAPDelta}{+0.069}
\providecommand{\PromptPrototypeDevelopmentFOneDelta}{+0.176}
\providecommand{\PromptPrototypeDevelopmentFOneDeltaLow}{+0.145}
\providecommand{\PromptPrototypeDevelopmentFOneDeltaHigh}{+0.210}
\providecommand{\PromptPrototypeMinimumFoldGain}{+0.146}
\providecommand{\FinalHoldoutImageCount}{100}
\providecommand{\FinalHoldoutGTCount}{169}
\providecommand{\SmallFinalPrecision}{0.645}
\providecommand{\SmallFinalRecall}{0.698}
\providecommand{\SmallFinalFOne}{0.670}
\providecommand{\SmallFinalTP}{118}
\providecommand{\SmallFinalFP}{65}
\providecommand{\SmallFinalFN}{51}
\providecommand{\SmallFinalAP}{0.559}
\providecommand{\SmallFinalAPFifty}{0.716}
\providecommand{\SmallFinalAPSeventyFive}{0.637}
\providecommand{\SmallTunedFinalPrecision}{0.689}
\providecommand{\SmallTunedFinalRecall}{0.657}
\providecommand{\SmallTunedFinalFOne}{0.673}
\providecommand{\MediumFinalPrecision}{0.752}
\providecommand{\MediumFinalRecall}{0.680}
\providecommand{\MediumFinalFOne}{0.714}
\providecommand{\MediumFinalTP}{115}
\providecommand{\MediumFinalFP}{38}
\providecommand{\MediumFinalFN}{54}
\providecommand{\MediumFinalAP}{0.605}
\providecommand{\MediumFinalAPFifty}{0.751}
\providecommand{\MediumFinalAPSeventyFive}{0.663}
\providecommand{\ModelFOneDelta}{+0.044}
\providecommand{\ModelFOneDeltaLow}{+0.008}
\providecommand{\ModelFOneDeltaHigh}{+0.082}
\providecommand{\ModelFOneDeltaVsSmallTuned}{+0.042}
\providecommand{\ModelFOneDeltaVsSmallTunedLow}{+0.007}
\providecommand{\ModelFOneDeltaVsSmallTunedHigh}{+0.076}
\providecommand{\ModelMAPDelta}{+0.046}
\providecommand{\MediumModelName}{yolov8m-worldv2.pt}
\providecommand{\MediumImageSize}{640}
\providecommand{\MediumCarThreshold}{0.35}
\providecommand{\MediumCouchThreshold}{0.50}
\providecommand{\MediumAirplaneThreshold}{0.30}
\providecommand{\MediumCupThreshold}{0.50}
\providecommand{\BaselinePrecision}{0.697}
\providecommand{\BaselineRecall}{0.664}
\providecommand{\BaselineFOne}{0.680}
\providecommand{\BaselineValidationFOne}{0.683}
\providecommand{\BaselineTP}{85}
\providecommand{\BaselineFP}{37}
\providecommand{\BaselineFN}{43}
\providecommand{\BaselineAP}{0.547}
\providecommand{\BaselineAPFifty}{0.727}
\providecommand{\BaselineAPSeventyFive}{0.539}
\providecommand{\BaselineARHundred}{0.690}
\providecommand{\BaselineAPSmall}{0.519}
\providecommand{\BaselineAPMedium}{0.467}
\providecommand{\BaselineAPLarge}{0.715}
\providecommand{\BaselineFOneCILow}{0.606}
\providecommand{\BaselineFOneCIHigh}{0.753}
\providecommand{\BlurPrecision}{0.857}
\providecommand{\BlurRecall}{0.281}
\providecommand{\BlurFOne}{0.424}
\providecommand{\BlurTP}{36}
\providecommand{\BlurFP}{6}
\providecommand{\BlurFN}{92}
\providecommand{\LowlightPrecision}{0.704}
\providecommand{\LowlightRecall}{0.594}
\providecommand{\LowlightFOne}{0.644}
\providecommand{\GammaPrecision}{0.714}
\providecommand{\GammaRecall}{0.625}
\providecommand{\GammaFOne}{0.667}
\providecommand{\GammaValue}{0.50}
\providecommand{\CanonicalPromptFOne}{0.680}
\providecommand{\SynonymPromptFOne}{0.599}
\providecommand{\HypernymPromptFOne}{0.436}
\providecommand{\CarCanonicalRecall}{0.632}
\providecommand{\CarSynonymRecall}{0.132}
\providecommand{\CarHypernymRecall}{0.000}
\providecommand{\CouchPhotoRecall}{0.846}
\providecommand{\CouchSceneRecall}{0.846}
\providecommand{\CupCanonicalRecall}{0.595}
\providecommand{\CupSynonymRecall}{0.514}
\providecommand{\PrimaryStrategy}{category\_subset\_nms}
\providecommand{\PrimaryValidationFOne}{0.764}
\providecommand{\PrimaryPrecision}{0.712}
\providecommand{\PrimaryRecall}{0.578}
\providecommand{\PrimaryFOne}{0.638}
\providecommand{\PrimaryAP}{0.502}
\providecommand{\PrimaryAPFifty}{0.706}
\providecommand{\PrimaryDelta}{-0.042}
\providecommand{\PrimaryDeltaLow}{-0.098}
\providecommand{\PrimaryDeltaHigh}{+0.017}
\providecommand{\NaiveValidationFOne}{0.661}
\providecommand{\NaivePrecision}{0.698}
\providecommand{\NaiveRecall}{0.688}
\providecommand{\NaiveFOne}{0.693}
\providecommand{\NaiveAP}{0.535}
\providecommand{\NaiveAPFifty}{0.736}
\providecommand{\NaiveDelta}{+0.013}
\providecommand{\NaiveDeltaLow}{-0.021}
\providecommand{\NaiveDeltaHigh}{+0.047}
\providecommand{\FixedValidationFOne}{0.719}
\providecommand{\FixedPrecision}{0.766}
\providecommand{\FixedRecall}{0.562}
\providecommand{\FixedFOne}{0.649}
\providecommand{\FixedAP}{0.507}
\providecommand{\FixedAPFifty}{0.682}
\providecommand{\FixedDelta}{-0.031}
\providecommand{\FixedDeltaLow}{-0.080}
\providecommand{\FixedDeltaHigh}{+0.014}
\providecommand{\CalibratedValidationFOne}{0.743}
\providecommand{\CalibratedPrecision}{0.798}
\providecommand{\CalibratedRecall}{0.523}
\providecommand{\CalibratedFOne}{0.632}
\providecommand{\CalibratedAP}{0.497}
\providecommand{\CalibratedAPFifty}{0.685}
\providecommand{\FusionValidationFOne}{0.717}
\providecommand{\FusionPrecision}{0.847}
\providecommand{\FusionRecall}{0.562}
\providecommand{\FusionFOne}{0.676}
\providecommand{\FusionAP}{0.516}
\providecommand{\FusionAPFifty}{0.713}
\providecommand{\GPUName}{NVIDIA GeForce RTX 4090}
\providecommand{\DriverVersion}{580.173.02}
\providecommand{\PythonVersion}{3.10.20}
\providecommand{\TorchVersion}{2.4.1+cu124}
\providecommand{\UltralyticsVersion}{8.4.112}
\providecommand{\FixedCarVariant}{photo\_template}
\providecommand{\SubsetCarVariants}{canonical + photo\_template}
\providecommand{\SubsetCarThreshold}{0.30}
\providecommand{\FixedCouchVariant}{canonical}
\providecommand{\SubsetCouchVariants}{plural}
\providecommand{\SubsetCouchThreshold}{0.20}
\providecommand{\FixedAirplaneVariant}{plural}
\providecommand{\SubsetAirplaneVariants}{scene\_template}
\providecommand{\SubsetAirplaneThreshold}{0.30}
\providecommand{\FixedCupVariant}{photo\_template}
\providecommand{\SubsetCupVariants}{hypernym + scene\_template}
\providecommand{\SubsetCupThreshold}{0.40}


\definecolor{Accent}{HTML}{275D8C}
\definecolor{DarkText}{HTML}{242A30}
\definecolor{Muted}{HTML}{59636E}
\definecolor{LightRule}{HTML}{D8DEE5}

\captionsetup{
  font=small,
  labelfont={bf,color=Accent},
  justification=centering
}
\setlist{nosep,leftmargin=2em}
\setlength{\parindent}{1em}
\setlength{\parskip}{0.10em}
\renewcommand{\arraystretch}{1.13}
\titlespacing*{\section}{0pt}{1.2ex plus 0.2ex minus 0.2ex}{0.5ex}
\titlespacing*{\subsection}{0pt}{0.9ex plus 0.2ex minus 0.2ex}{0.35ex}
\pagestyle{fancy}
\fancyhf{}
\lhead{\small YOLO-World Failure Case Analysis}
\rhead{\small 映像メディア学レポート}
\cfoot{\thepage}
\renewcommand{\headrulewidth}{0.3pt}

\newcommand{\result}[1]{\textbf{\color{Accent}#1}}
\newcommand{\code}[1]{\texttt{#1}}

\begin{document}

\begin{center}
  {\LARGE\sffamily\bfseries YOLO-WorldのFailure Case Analysisと改善\par}
  \vspace{0.35em}
  {\large\sffamily blur-aware前処理とcanonical-anchored text prototype\par}
\end{center}
\vspace{0.6em}

\noindent
\begin{tabularx}{\linewidth}{@{}>{\bfseries}lX>{\bfseries}lX@{}}
氏名 & 杉中史哉 & 学籍番号 & 48-256421 \\
研究科 & 情報理工学系研究科 & 学年 & 修士2年 \\
専攻 & 電子情報学専攻 & 所属研究室 & 菅野研究室 \\
研究テーマ & \multicolumn{3}{l}{競技用義足ユーザーを対象とした姿勢推定および義足ブレードの3次元形状再構築} \\
\end{tabularx}
\vspace{0.5em}
\hrule
\vspace{0.7em}

\section*{概要}

CVPR 2024のopen-vocabulary物体検出器YOLO-WorldをCOCO 2017 valで評価した。
強いGaussian blurとprompt表現の変更でRecall/F1が低下するFailure Caseを確認し、
入力補正による改善を検証した。独立評価用100枚でblur-aware Wiener前処理は
blur 2条件の平均F1を\SameModelBaselineBlurMeanFOne から
\SameModelImprovedBlurMeanFOne（差\result{\SameModelBlurMeanFOneDelta}、
95\%区間$[\SameModelBlurMeanFOneDeltaLow,\SameModelBlurMeanFOneDeltaHigh]$）へ改善した。
別の独立評価用100枚でcanonical-anchored text prototypeはsynonym/hypernym平均F1を
\PromptPrototypeBaselineTargetMeanFOne から\PromptPrototypeImprovedTargetMeanFOne へ
改善した。

\section{対象論文概要}

対象論文はChengらの``YOLO-World: Real-Time Open-Vocabulary Object Detection''
\cite{yoloworld}であり、CVPR 2024に採択された。学習時に固定されたカテゴリだけを
扱う通常の物体検出器と異なり、テキストで指定した対象を実時間で検出する
open-vocabulary物体検出器である。

\section{手法理解}

YOLO-WorldはYOLOv8 detector、CLIP text encoder、RepVL-PANから構成される。
画像encoderが多段階の画像特徴を抽出し、text encoderがカテゴリ名をembeddingへ変換する。
RepVL-PANではText-guided CSPLayerが言語情報を画像特徴へ注入し、
Image-Pooling Attentionが画像情報を使ってtext embeddingを強化する。
検出headはboxとobject embeddingを予測し、object--text類似度をカテゴリscoreとして扱う。

学習では正例・負例の語彙を含むonline vocabularyとregion--text contrastive lossを用い、
画像領域と語彙を対応付ける。推論ではpromptを先にoffline vocabulary embeddingへ変換する
prompt-then-detect方式を採る。このembeddingは検出層の重みへ再パラメータ化できるため、
入力画像ごとにtext encoderを再実行せず、open-vocabulary検出を効率化できる。

\section{実行環境}

検出器にはUltralytics版の公開重み
\code{\SameModelName}を用い、ローカルのGPU \GPUName で実行した。
driverは\DriverVersion、Pythonは\PythonVersion、PyTorchは\TorchVersion、
Ultralyticsは\UltralyticsVersion である。評価では画像・GT box・検出器をそろえ、
画像劣化とprompt表現の影響、および入力補正の効果を比較した。

\subsection{再現手順}

データ抽出、推論、評価、図表生成では共通の設定と乱数seedを用いた。
ソースコード、実験条件、実行手順は
\href{https://github.com/fumisugi/visual_media}{GitHub repository}に記載した。

\section{Failure Case Analysis}

\subsection{評価設計と指標}

対象カテゴリは\code{car / couch / airplane / cup}である。COCO annotationのうち
crowdを除き、bbox面積が1,024 pixel以上の対象を含む画像を固定seedで抽出した。
F1評価ではconfidence降順に予測を処理し、同一カテゴリかつIoU 0.50以上の未対応GTへ
greedy matchingした。一つのGTに複数予測が重なる場合、最初の一つだけをTrue Positive、
残りをFalse Positiveとした。

データは互いに重ならない5組、計500枚を用いた。最初の100枚はFailure Caseとprompt
比較に用い、validation 30枚で条件を選択し、test 70枚で評価した。
blur前処理は\SameModelDevelopmentImageCount 枚で方式とパラメータを選び、
非重複の\SameModelHoldoutImageCount 枚で評価した。prompt prototypeは400枚で
混合係数を選び、別の非重複100枚で評価した。評価用データを方式選択には使用していない。
適格な\code{airplane}画像が限られたため、blur評価セットのカテゴリ構成は
\code{car/couch/airplane/cup}=28/28/16/28とした。

\begin{table}[H]
\centering
\caption{Canonical baselineのtest指標}
\label{tab:baseline}
\begin{tabular}{lrrrrrrrr}
\toprule
P & R & F1 & mAP & AP50 & AP75 & AR@100 & F1 95\% CI \\
\midrule
\BaselinePrecision &
\BaselineRecall &
\BaselineFOne &
\BaselineAP &
\BaselineAPFifty &
\BaselineAPSeventyFive &
\BaselineARHundred &
$[\BaselineFOneCILow,\BaselineFOneCIHigh]$ \\
\bottomrule
\end{tabular}
\end{table}

COCO mAPはIoU 0.50から0.95まで0.05刻みで平均したAP、AP50/AP75は各IoUにおける
interpolated Precision--Recall曲線の面積、AR@100は1画像・カテゴリ当たり最大100予測の
平均Recallである。APはraw confidence 0.01以上のranked predictionsから算出し、F1は
固定またはvalidationで選択した動作thresholdで算出した。不確実性はtest画像を単位に
2,000回再標本化したbootstrapで評価した。

\subsection{画像劣化}

入力画像、prompt、GT boxを固定し、明るさ倍率
$1.0/0.5/0.25$とGaussian blur $\sigma=0/2/4$を比較した。
通常画像からblur $\sigma=4$へ変えるとTPは
\BaselineTP 件から\BlurTP 件、FNは\BaselineFN 件から\BlurFN 件へ変化した。
Precisionは\BaselinePrecision から\BlurPrecision へ上がった一方、
Recallは\BaselineRecall から\BlurRecall へ低下した。これは低confidenceの予測自体が
減ったためであり、誤ったboxが急増するというより、対象に対応するboxを出せなくなる。

\begin{figure}[htbp]
\centering
\includegraphics[width=0.80\linewidth]{failure_image_corruptions.png}
\caption{画像劣化条件ごとのPrecision、Recall、F1。
同じtest \TestImageCount 枚、confidence 0.25、matching IoU 0.50で比較。}
\label{fig:degradation}
\end{figure}

\begin{figure}[htbp]
\centering
\begin{subfigure}[t]{0.45\linewidth}
  \centering
  \includegraphics[width=\linewidth]{baseline_success.png}
  \caption{通常画像ではGTと予測が一致。}
\end{subfigure}\hfill
\begin{subfigure}[t]{0.45\linewidth}
  \centering
  \includegraphics[width=\linewidth]{failure_under_gaussian_blur_sigma4.png}
  \caption{blur $\sigma=4$では対象を検出できない。}
\end{subfigure}
\caption{画像劣化Failure Caseの代表例。緑がGT、赤が予測。}
\label{fig:blur-example}
\end{figure}

ぼけによって輪郭や局所textureが失われ、候補boxへ十分なconfidenceを与えられなくなった
可能性がある。ただし内部特徴を可視化していないため、これは結果に整合する仮説であり、
因果的に確認した説明ではない。

\subsection{Prompt表現}

元画像を固定し、canonical、synonym、hypernym、写真template
（``a photo of a \dots''）、scene template（``a \dots in the image''）、
複数形の6種類を比較した。全カテゴリ合計ではcanonical、synonym、hypernymのF1が
\CanonicalPromptFOne、\SynonymPromptFOne、\HypernymPromptFOne となった。

\begin{figure}[H]
\centering
\begin{subfigure}[t]{0.48\linewidth}
  \centering
  \includegraphics[width=\linewidth]{failure_prompt_wording.png}
  \caption{基本3表現のP/R/F1。}
\end{subfigure}\hfill
\begin{subfigure}[t]{0.50\linewidth}
  \centering
  \includegraphics[width=\linewidth]{prompt_variant_recall_heatmap.png}
  \caption{6表現のカテゴリ別Recall。}
\end{subfigure}
\caption{Prompt wordingによる検出性能の変化。画像とGTは固定。}
\label{fig:prompt-sensitivity}
\end{figure}

\code{car}のtest Recallはcanonical \CarCanonicalRecall、synonym
\CarSynonymRecall、hypernym \CarHypernymRecall だった。一方、
\code{couch}のphoto/scene templateは\CouchPhotoRecall/\CouchSceneRecall、
\code{cup}のcanonical/synonymは\CupCanonicalRecall/\CupSynonymRecall だった。
同義語や説明句の効果は一様ではなく、カテゴリごとに適切な表現が異なる。

この差は、text encoderが辞書的な同義関係だけでなく、事前学習データ中の語と画像の
共起を反映するためと考えられる。広い上位概念は対象領域との対応が曖昧になり、
canonical termよりconfidenceが下がる場合がある。したがって、全カテゴリへ同じ種類の
言い換えを一律適用することは安全ではない。

\section{改善内容}

\subsection{Canonical-anchored text prototype}

最初にnaive NMS ensemble、カテゴリ別prompt/threshold、prompt subset、weighted box
fusionをvalidationだけで調整した。しかしvalidation全体F1で選ばれた
\code{\PrimaryStrategy}はtest F1を\BaselineFOne から\PrimaryFOne へ低下させ、
差\PrimaryDelta の95\%区間$[\PrimaryDeltaLow,\PrimaryDeltaHigh]$も0を含んだ。
naive 3-prompt NMSも差\NaiveDelta、区間
$[\NaiveDeltaLow,\NaiveDeltaHigh]$で、mAPは\BaselineAP から\NaiveAP へ下がった。
複数promptのboxを足す方法ではFP増加とvalidation過適合を避けられなかった。

そこで検出後のboxではなくtext embeddingを統合した。各カテゴリについて入力語
$e_{\mathrm{input}}$とcanonical語$e_{\mathrm{canon}}$をCLIPで単位ベクトル化し、
\[
e_{\mathrm{proto}} =
\operatorname{normalize}\left(
(1-\alpha)e_{\mathrm{input}}+\alpha e_{\mathrm{canon}}
\right)
\]
を1クラス1個のprototypeとしてYOLO-Worldへ与える。canonical入力では元のembeddingを
そのまま使い、synonym/hypernymだけを補正する。開発用
\PromptPrototypeDevelopmentImageCount 枚を4 foldに分け、
$\alpha\in\{0.25,0.50,0.75\}$を比較した結果、
$\alpha=\PromptPrototypeAnchorWeight$を選択した。
synonym/hypernym平均F1差は\PromptPrototypeDevelopmentFOneDelta、95\%区間
$[\PromptPrototypeDevelopmentFOneDeltaLow,\PromptPrototypeDevelopmentFOneDeltaHigh]$、
最小fold差も\PromptPrototypeMinimumFoldGain だった。

\subsection{Blur-aware Wiener前処理}

改善前後でモデル、prompt、入力解像度、confidence thresholdを固定し、入力前処理だけを
変更した。まず開発用30枚で古典的なunsharp maskとWiener deconvolutionを比較し、
次に互いに重ならない3組、計\SameModelDevelopmentImageCount 枚を3 foldとして
候補を比較した。各画像の離散Laplacian varianceからblurを判定し、blurの場合だけ
Wiener deconvolution（regularization \WienerRegularization、元画像とのblend
\WienerBlend）を適用した。開発用300枚ではblur 2条件の平均F1差が
\SameModelDevelopmentFOneDelta、paired bootstrap 95\%区間は
$[\SameModelDevelopmentFOneDeltaLow,\SameModelDevelopmentFOneDeltaHigh]$で、
最小fold差も\SameModelMinimumFoldGain だった。

\subsection{Gamma correction}

明るさ25\%へ人工的に暗化した画像にpower-law gamma correctionを適用した。
validationで$\gamma=0.40/0.50/0.60/0.70$を比較した。
\ifnum\GammaValidationTieCount>1
\GammaValidationTieCount 候補が最高F1で同率だったため、同率候補の中から
変形が最も弱く1.0に近い$\gamma=\GammaValue$をtie-breakで選択した。
\else
最高F1だった$\gamma=\GammaValue$を選択した。
\fi

\section{改善前後比較}

\subsection{Prompt prototype}

\begin{figure}[H]
\centering
\includegraphics[width=0.92\linewidth]{prompt_prototype_final_holdout.png}
\caption{評価用COCO val \PromptPrototypeHoldoutImageCount 枚におけるprompt表現別比較。
canonicalは不変で、synonymとhypernymのF1/mAPが上昇した。}
\label{fig:prompt-prototype}
\end{figure}

開発用400枚で選択した混合係数を、非重複の評価用100枚
（GT \PromptPrototypeHoldoutGTCount 個）に適用した。適格なairplane画像が
5枚に限られたため、最小box面積256、カテゴリ構成31/32/5/32とした。
synonym/hypernym平均F1は
\PromptPrototypeBaselineTargetMeanFOne から\PromptPrototypeImprovedTargetMeanFOne
（差\result{\PromptPrototypeTargetMeanFOneDelta}、95\%区間
$[\PromptPrototypeTargetMeanFOneDeltaLow,\PromptPrototypeTargetMeanFOneDeltaHigh]$）、
平均mAPは\PromptPrototypeBaselineTargetMeanAP から
\PromptPrototypeImprovedTargetMeanAP（差\PromptPrototypeTargetMeanAPDelta）へ上昇した。
条件別F1はsynonymで\PromptPrototypeBaselineSynonymFOne から
\PromptPrototypeImprovedSynonymFOne、hypernymで
\PromptPrototypeBaselineHypernymFOne から\PromptPrototypeImprovedHypernymFOne だった。
canonical F1は\PromptPrototypeBaselineCanonicalFOne のまま変わらない。

\begin{figure}[H]
\centering
\begin{subfigure}[t]{0.48\linewidth}
  \centering
  \includegraphics[width=\linewidth]{prompt_prototype_baseline.png}
  \caption{hypernym入力では対象を検出できなかった。}
\end{subfigure}\hfill
\begin{subfigure}[t]{0.48\linewidth}
  \centering
  \includegraphics[width=\linewidth]{prompt_prototype_improvement.png}
  \caption{prototype補正後は2個の\code{car}を検出した。}
\end{subfigure}
\caption{Prompt prototypeによる改善例。緑がGT、赤が予測。}
\label{fig:prompt-prototype-example}
\end{figure}

\subsection{Blur-aware Wiener前処理}

\begin{figure}[H]
\centering
\includegraphics[width=0.90\linewidth]{same_model_blur_final_holdout.png}
\caption{評価用COCO val \SameModelHoldoutImageCount 枚におけるWiener前処理の比較。
cleanはほぼ維持し、強いblur（$\sigma=4$）を大きく回復した。}
\label{fig:same-model-blur}
\end{figure}

開発用300枚で決定した前処理を、非重複の評価用\SameModelHoldoutImageCount 枚
（car/couch/airplane/cup = 28/28/16/28、GT \SameModelHoldoutGTCount 個）に
適用した。blur 2条件の平均F1は
\SameModelBaselineBlurMeanFOne から\SameModelImprovedBlurMeanFOne へ上がり、
差は\result{\SameModelBlurMeanFOneDelta}、95\%区間
$[\SameModelBlurMeanFOneDeltaLow,\SameModelBlurMeanFOneDeltaHigh]$で、事前基準
（差$\geq0.05$かつ区間下限$>0$）を満たした。clean F1は
\SameModelBaselineCleanFOne から\SameModelImprovedCleanFOne
（差\SameModelCleanFOneDelta）で、許容低下0.01以内だった。

改善は主に$\sigma=4$で生じ、F1は
\SameModelBaselineBlurFourFOne から\SameModelImprovedBlurFourFOne、mAPは
\SameModelBaselineBlurFourAP から\SameModelImprovedBlurFourAP、Recallは
\SameModelBaselineBlurFourRecall から\SameModelImprovedBlurFourRecall へ上昇した。
一方$\sigma=2$ではF1
\SameModelBaselineBlurTwoFOne から\SameModelImprovedBlurTwoFOne とほぼ同じで、
mAPは\SameModelBaselineBlurTwoAP から\SameModelImprovedBlurTwoAP へわずかに低下した。
したがって「全blurに一様に有効」ではなく「強いGaussian blurを回復する改善」である。

\begin{figure}[htbp]
\centering
\begin{subfigure}[t]{0.48\linewidth}
  \centering
  \includegraphics[width=\linewidth]{same_model_blur_baseline.png}
  \caption{blur $\sigma=4$では2機を見落とした。}
\end{subfigure}\hfill
\begin{subfigure}[t]{0.48\linewidth}
  \centering
  \includegraphics[width=\linewidth]{same_model_blur_improvement.png}
  \caption{Wiener前処理後は同じモデルが2機を検出した。}
\end{subfigure}
\caption{同一画像・同一検出器における改善例。緑がGT、赤が予測。}
\label{fig:same-model-example}
\end{figure}

blur状態は\SameModelDetectorCorrect/\SameModelDetectorTotal 入力で正しく分類された。
clean 100枚中\SameModelCleanCorrect 枚をcleanと判定し、1枚だけ誤って補正したため、
cleanでF1がわずかに低下した。

\subsection{Gamma correction}

test F1は
\LowlightFOne から\GammaFOne、Recallは\LowlightRecall から\GammaRecall へ変化した。

効果は小さく、単一splitで選ばれた$\gamma$の安定性は未確認である。また人工的な暗化は
実環境のshot noiseやsensor noiseを再現せず、gamma correctionによるnoise増幅も
評価していない。したがって、実写low-lightへの一般化は未確認である。

\section{Limitations}

\begin{itemize}
  \item 互いに重ならない計500枚を使ったが、対象は4カテゴリに限られ、COCO全体や
  別datasetの性能を表さない。適格なairplane画像が少なく、blur評価セットでは16枚、
  prompt評価セットでは5枚であり、カテゴリ構成は完全には均衡していない。
  \item text prototypeは入力語が既知canonicalカテゴリの同義・上位概念だと対応付けられる
  前提を持つ。未知の自由記述、複数概念、誤ったmappingへの頑健性は未評価である。
  \item blur-aware前処理は独立100枚で確認したが、利得は$\sigma=4$へ集中し、
  $\sigma=2$のmAPはわずかに低下し、clean 1枚を誤補正した。追加時間も未測定である。
  \item 人工的劣化は実環境のnoiseやmotion blurを完全には再現しない。内部特徴も
  可視化していないため、失敗原因の説明は診断仮説に留まる。
\end{itemize}

\section{生成AI利用内容}

著者が対象論文を選定し、blurとprompt表現をFailure Caseとして設定した。
また、評価画像数、評価指標、改善の採否基準、結果の解釈とLimitationsを決定した。
OpenAI Codexは、著者の指示に基づくPython/LaTeX実装、実験実行、集計、図表作成、
整合性確認の補助に使用した。生成されたコードと文章は実行結果と照合し、著者が
主張範囲、考察、レポート表現を確認・修正した。
詳細な利用記録は\href{https://github.com/fumisugi/visual_media/blob/main/doc/AI_USAGE_LOG.md}
{GitHub上の生成AI利用ログ}に記載した。

\section{考察}

YOLO-Worldは強いblurとprompt変更でFailure Caseを示した。3-prompt NMSや
validation選択法は一般化しなかったが、boxを増やさずtext embeddingをcanonicalへ
anchorする方法は、独立評価用100枚でsynonym/hypernym平均F1を
\PromptPrototypeBaselineTargetMeanFOne から\PromptPrototypeImprovedTargetMeanFOne へ
改善し、mAPも上昇した。canonical入力の結果は変えていない。

blur-aware Wiener前処理は開発用300枚で方式を選択し、非重複の評価用100枚で
検証した。比較条件をそろえた結果、
blur 2条件の平均F1を\SameModelBaselineBlurMeanFOne から
\SameModelImprovedBlurMeanFOne へ改善し、paired区間の下限も0より上だった。
特に強いblurでF1とmAPがともに回復した一方、軽いblurとcleanでの小さな不利益も示した。
2つのFailure Caseそれぞれへ直接対応し、手法選択用データと評価用データを分離したことが
本実験の主要な知見である。

\footnotesize
\begin{thebibliography}{9}
\setlength{\itemsep}{0pt}
\setlength{\parskip}{0pt}
\bibitem{yoloworld}
T. Cheng et al.,
``YOLO-World: Real-Time Open-Vocabulary Object Detection,''
CVPR, 2024.
\href{https://openaccess.thecvf.com/content/CVPR2024/html/Cheng_YOLO-World_Real-Time_Open-Vocabulary_Object_Detection_CVPR_2024_paper.html}{CVF Open Access}

\bibitem{official}
AILab-CVC,
``Official YOLO-World implementation.''
\href{https://github.com/AILab-CVC/YOLO-World}{GitHub repository}

\bibitem{ultralytics}
Ultralytics,
``YOLO-World documentation.''
\href{https://docs.ultralytics.com/models/yolo-world/}{Ultralytics documentation}

\bibitem{coco}
COCO Consortium,
``Common Objects in Context.''
\href{https://cocodataset.org/}{COCO website}
\end{thebibliography}

\end{document}
